\chapter*{第五部分收束：把复杂工具带回可信交付}
\addcontentsline{toc}{chapter}{第五部分收束：把复杂工具带回可信交付}

Part V 让我们接触了更复杂的模型和更强的 AI 工具。到这里，一个学生已经可能写出神经网络训练代码、构造卷积特征、解释自编码器重构误差、描述注意力机制，也能用 LLM 协助调试、整理文献和搭建工具调用流程。

但本部分真正要留下的并不是“我会用复杂工具”，而是“我知道复杂工具应该怎样被纳入可信交付”。这两者差别很大。

\section*{复杂度必须有理由}

进入 Capstone 项目前，任何复杂方法都应能回答三个问题：

\begin{enumerate}
    \item 它解决了简单 baseline 的哪个具体不足；
    \item 它引入了哪些新的风险和验证成本；
    \item 它的输出如何回到科学图表、错误样本和报告结论。
\end{enumerate}

如果一个深度模型只是在分数上略有提升，却让数据需求、解释成本和复现难度大幅上升，那么它未必是更好的课程项目选择。相反，一个简单模型如果能被完整解释、稳定复现、清楚展示限制，往往更适合本科 capstone。

\section*{AI 助手也要有交付边界}

LLM 和 agentic workflow 可以显著提高效率，但它们在课程项目中必须留下可审计痕迹：

\begin{itemize}
    \item 哪些代码由 AI 起草，哪些由你验证；
    \item 哪些文献摘要来自 AI 辅助，哪些已经回到原文核对；
    \item 哪些工具调用成功，哪些失败或被人工接管；
    \item 哪些内容需要引用、署名或版权检查；
    \item 最终报告中哪些判断是你自己承担责任的科学解释。
\end{itemize}

这不是为了限制 AI，而是为了让 AI 真正成为科研助手，而不是不可审计的黑箱作者。

\section*{进入 Capstone 前的选择原则}

Part VI 会要求学生完成一个完整项目。此时最重要的能力不是会多少模型，而是会做选择：

\begin{itemize}
    \item 能用简单方法回答的问题，不要为了炫技而复杂化；
    \item 需要复杂模型的问题，必须先有简单 baseline；
    \item 需要 AI 助手参与的流程，必须保留验证记录；
    \item 所有最终结论都必须能回到数据、代码、图表和限制说明。
\end{itemize}

换句话说，现代 AI 的正确位置不是替代科学判断，而是帮助你更快、更稳地抵达需要判断的地方。
